\documentclass[a0paper,portrait, blockverticalspace=4em, colspace=2em]{tikzposter}
\tikzposterlatexaffectionproofoff
\usepackage{graphicx}
\usepackage{amsmath}
\usepackage{amssymb}
\usepackage{hyperref}
\usepackage{multicol}
\usepackage{booktabs}
\usetikzlibrary{calc}

% ---------------------------------------------------------------------------
%  Colours
% ---------------------------------------------------------------------------
\definecolor{maincolor}{HTML}{9e2b2f}
\definecolor{secondarycolor}{RGB}{158, 43, 47}
\definecolor{lightred}{RGB}{237, 199, 201}

\definecolor{tid}{RGB}{254,128,2}   % orange  — Total Ionizing Dose
\definecolor{niel}{RGB}{0,140,0}    % green   — Non-Ionizing Energy Loss
\definecolor{see}{RGB}{190,0,62}    % red     — Single Event Effects
\definecolor{sel}{RGB}{0,90,170}    % blue    — Single Event Latch-up

\newcommand{\looseitems}{\normalsize\begin{itemize}\setlength{\itemsep}{0.3em}}
\newcommand{\tightitems}{\normalsize\begin{itemize}\setlength{\itemsep}{0.15em}}

\newcommand{\captiontext}[1]{%
    {\small\renewcommand{\baselinestretch}{0.88}\selectfont#1}}

% ---------------------------------------------------------------------------
%  Theme
% ---------------------------------------------------------------------------
\usetheme{Default}
\definecolorstyle{atlascolors}{
    \colorlet{backgroundcolor}{white}
    \colorlet{titlefgcolor}{white}
    \colorlet{titlebgcolor}{maincolor}
    \colorlet{blocktitlefgcolor}{white}
    \colorlet{blocktitlebgcolor}{maincolor}
    \colorlet{blockbodyfgcolor}{black}
    \colorlet{blockbodybgcolor}{lightred!25}
}{}
\usecolorstyle{atlascolors}

% Title bar with ATLAS logo at right
\definetitlestyle{sampletitle}{width=760mm, roundedcorners=20, linewidth=2pt,
  innersep=10pt, titletotopverticalspace=6mm, titletoblockverticalspace=8mm}{%
  \begin{scope}[line width=\titlelinewidth, rounded corners=\titleroundedcorners]
    \draw[color=blocktitlebgcolor, fill=titlebgcolor]
      (\titleposleft,\titleposbottom) rectangle (\titleposright,\titlepostop);
  \end{scope}
  \node[anchor=east, fill=white, rounded corners=10pt, inner sep=10pt, xshift=-22mm]
    at ($(\titleposright,\titlepostop)!0.5!(\titleposright,\titleposbottom)$)
    {
\includegraphics[height=8.6cm]{logos/atlas_transparent.png}};}
\usetitlestyle{sampletitle}

\title{\parbox{0.74\linewidth}{\centering\huge
    Irradiation Studies and Design Optimization of the\\[2mm]ATLAS Tile Calorimeter for the High-Luminosity LHC}}
\author{Eduardo Valdes Santurio \quad on behalf of the ATLAS Tile Calorimeter System}
\institute{Stockholm University}
\date{}

\setlength{\columnsep}{0.5cm}

\begin{document}
\maketitle

\begin{columns}

% ===========================================================================
%  COLUMN 1
% ===========================================================================
\column{0.333}

\block{The ATLAS Tile Calorimeter Phase-II Upgrade}{
    \tightitems
        \item As the primary hadronic sampling calorimeter of ATLAS, TileCal is crucial for
        measuring hadronic jets, missing transverse energy, hadronically decaying $\tau$-leptons,
        and for muon identification.
        \item To cope with the unprecedented instantaneous luminosity and pile-up of the
        \textbf{High-Luminosity LHC}, the on-detector electronics system is being \textbf{fully replaced}.
        \item Upgraded on-detector electronics will digitize PMT signals at $40\,$MHz and transmit
        the continuous data stream to the off-detector systems via high-speed optical links.
    \end{itemize}
    \vspace{0.12cm}
    \begin{center}
        \includegraphics[width=1.0\linewidth]{assets/fig_minidrawer_blockdiagram.png}
    \end{center}
    \vspace{-0.2cm}
    \captiontext{Block diagram of the Minidrawer readout architecture for HL-LHC.}
}

\block{Radiation Environment \& Requirements}{
    \tightitems
        \item The HL-LHC will deliver an integrated luminosity of 4000\,fb$^{-1}$. The new on-detector
        hardware faces a vastly harsher radiation environment than originally designed for.
        \item Every component must be rigorously qualified against \textbf{\textcolor{tid}{Total Ionizing Dose (TID)}},
        \textbf{\textcolor{niel}{Non-Ionizing Energy Loss (NIEL)}}, and \textbf{\textcolor{see}{Single Event Effects (SEE/SEL)}}.
        \item Radiation criteria are based on ATLAS simulations for the most-exposed regions,
        scaled by rigorous safety factors (SF) accounting for simulation uncertainty (1.5),
        low-dose-rate effects (3--5), and non-monoenergetic beams (1.3).
        \item Final Daughterboard (DB) qualification targets (with all SFs applied):
        \textbf{\textcolor{tid}{TID $\approx$ 108\,Gy}},
        \textbf{\textcolor{niel}{NIEL $\approx 13.16\times10^{12}$\,n/cm$^2$}},
        \textbf{\textcolor{see}{SEE $\approx 4.54\times10^{11}$\,p/cm$^2$}}.
    \end{itemize}
}

\block{The Link \& Control Daughterboard (DB6)}{
    \tightitems
        \item The DB is the central communication hub of the Minidrawer, responsible for high-speed
        data transmission to the off-detector Tile PreProcessor, and distributing configuration and clocks.
        \item It achieves exceptional radiation tolerance via a \textbf{double-redundant} A/B architecture
        enforced by \textbf{Triple Modular Redundancy (TMR)}, \textbf{Xilinx Soft Error Mitigation (SEM)},
        \textbf{GBT-FEC} error correction, and \textcolor{sel}{SEL}-protection circuitry.
        \item Core components include 4$\times$ commercial \textbf{SFP+} transceivers, custom CERN
        \textbf{GBTx} ASICs, Xilinx \textbf{Kintex UltraScale (KU)}, and Microsemi \textbf{ProASIC3} FPGAs.
    \end{itemize}
    \vspace{0.12cm}
    \begin{center}
        \includegraphics[width=1.0\linewidth]{assets/fig_db6_blockdiagram.png}
    \end{center}
    \vspace{-0.2cm}
    \captiontext{DB6 block diagram: two redundant halves, KU + ProASIC FPGAs, GBTx and four SFP$+$ optical links.}
}

\block{References}{
    \small
    \begin{enumerate}\setlength{\itemsep}{0.2em}
        \item ATLAS Collaboration, \emph{Technical Design Report for the Phase-II Upgrade of the ATLAS Tile Calorimeter}, CERN-LHCC-2017-019 (2018).
        \item E. Valdes Santurio \textit{et al.}, \emph{Radiation studies performed on the High-Luminosity ATLAS TileCal link Daughterboard}, JINST \textbf{18} (2023) C04011.
        \item E. Valdes Santurio \textit{et al.}, \emph{Radiation studies on the ATLAS Link and Control Daughterboard for the HL-LHC Upgrade}, IEEE NSS MIC RTSD (2025).
    \end{enumerate}
}

% ===========================================================================
%  COLUMN 2
% ===========================================================================
\column{0.333}

\block{Total Ionizing Dose (TID) Testing}{
    \tightitems
        \item DB6 boards were irradiated at the CERN CC60 facility ($^{60}$Co source) to 220\,Gy
        at a dose rate of 3.37\,Gy/h.
        \item Additional proton irradiation was performed at the IFJ PAN facility in Kraków
        with 54\,MeV protons, reaching $\sim$400\,Gy at 2.88\,Gy/h.
        \item Temperatures and currents remained completely stable throughout the testing period,
        confirming that the DB6 sustains no significant \textcolor{tid}{TID} damage.
    \end{itemize}
    \vspace{0cm}
    \begin{center}
        \includegraphics[width=1.0\linewidth]{assets/fig_db6_tid_test.png}
    \end{center}
    \vspace{-0.3cm}
    \captiontext{DB6-3 board currents and FPGA temperature, stable over the full
    $\sim$400\,Gy proton \textcolor{tid}{TID} run (both halves).}
}

\block{Non-Ionizing Energy Loss (NIEL) Testing}{
    \tightitems
        \item A DB6 board (DB6-0) alongside test boards was irradiated at the TRIGA reactor
        (JSI Ljubljana) to a fluence of \textbf{$14\times10^{12}$\,n$_\text{eq}$/cm$^2$} (1\,MeV equivalent).
        \item The KU FPGAs demonstrated strong tolerance up to the exposed fluence with no degradation.
        \item The ProASIC FPGAs and IS25LP128F FLASH memories exhibited some after-effects where
        re-configuration failed initially. However, full functionality and re-programming capabilities
        were successfully recovered after annealing at room temperature.
    \end{itemize}
}

\block{Single Event Effects (SEE \& SEL)}{
    \tightitems
        \item \textcolor{sel}{SEL} testing was performed using 54\,MeV protons (Kraków) to a fluence
        of $2.62\times10^{11}$\,p/cm$^2$.
        \item Further tests with 226\,MeV protons irradiated Trenz TE0841 boards (using the same KU FPGA)
        at Mass General Hospital. The setup identified an unexpected \textcolor{sel}{SEL} risk 
        (cross section $\sim 1\times 10^{-10}$\,cm$^2$/device) in the Kintex UltraScale+ model initially 
        considered. 
        \item The testing setup demonstrated high reliability, leading to the selection of the standard
        KU architecture for production.
        \item Across all high-energy proton campaigns, \textbf{zero single-event latch-ups} were observed
        in any of the KU FPGAs, validating the intrinsic \textcolor{see}{SEE} robustness of the design.
    \end{itemize}
    \vspace{0.12cm}
    \begin{center}
        \includegraphics[width=1.0\linewidth]{assets/fig_system_tests_charm.jpeg}
    \end{center}
    \vspace{-0.2cm}
    \captiontext{System-level testing of the Daughterboard prototype at the CHARM facility.}
}

% ===========================================================================
%  COLUMN 3
% ===========================================================================
\column{0.333}

\block{Component Qualification Matrix}{
    \normalsize
    The DB6 incorporates commercial-off-the-shelf (COTS) parts requiring rigorous validation.
    The table summarizes key radiation testing results for these components:
    
    \vspace{0.3cm}
    \begin{center}
    \renewcommand{\arraystretch}{1.2}
    \begin{tabular}{llp{9cm}}
    \toprule
    \textbf{Component} & \textbf{Part Number} & \textbf{Qualification Result} \\
    \midrule
    SFP+ (AVAGO) & AFBR-709SMZ & \textbf{Failed} completely at $9\times 10^{12}$\,n/cm$^2$ \\
    SFP+ (CORETEK) & CT-A000NPP & \textbf{1/9 failed} at $9\times 10^{12}$\,n/cm$^2$ \\
    SFP+ (Direktronik) & 33-4248 & \textbf{Passed} up to $14\times 10^{12}$\,n/cm$^2$ \\
    SFP+ (FS-Europe) & SFP-10GSR-85 & \textbf{Passed} up to $14\times 10^{12}$\,n/cm$^2$ \\
    MOSFET & FDFMA3N109 & \textbf{Failed}: Shifted threshold voltage at 195\,Gy \\
    Oscillator (QUAD) & LTC6902IMS & \textbf{Failed}: 4/15 failed post-irradiation \\
    Schottky Diode & DAN217T146CT & \textbf{Passed} post-radiation tests \\
    Oscillator & 540BAA100 & \textbf{Passed} post-radiation tests \\
    Inst. Amp & INA333AIDRG & \textbf{Passed} post-radiation tests \\
    \bottomrule
    \end{tabular}
    \captiontext{Summary of DB6 COTS component irradiation tests.}
    \end{center}
}

  \block{SFP+ Optical Transceiver Matrix}{
      \tightitems
          \item Reliable, high-bandwidth data transmission at 9.6\,Gbps is absolutely critical for the readout system. 
          \item SFP+ components were evaluated under severe \textcolor{niel}{NIEL} conditions. Out of four tested variants, only two passed the mandatory $14\times10^{12}$\,n/cm$^2$ target.
          \item Additional \textcolor{tid}{TID} and \textcolor{sel}{SEL} testing needs to be carried out on the qualified SFP+ variants before their final integration into the DB6 production design.
      \end{itemize}
      \vspace{0.3cm}
      \begin{center}
      \small
      \begin{tabular}{l l c}
      \hline\hline
      \textbf{SFP+ Model} & \textbf{TID / NIEL Exposure Limit} & \textbf{Status} \\
      \hline
      AVAGO AFBR-709SMZ & Failed during early irradiation & \textcolor{red}{Rejected} \\
      CORETEK CT-A000NPP & 1/9 units failed @ $9\times 10^{12}$ & \textcolor{red}{Rejected} \\
      Direktronik 33-4248 & Passed 21$\times10^{12}$\,n/cm$^2$, $>$200\,Gy & \textcolor{niel}{Qualified} \\
      FS-Europe SFP-10GSR & Passed 21$\times10^{12}$\,n/cm$^2$, $>$200\,Gy & \textcolor{niel}{Qualified} \\
      \hline\hline
      \end{tabular}
      \end{center}
  }

\block{Design Optimizations \& Modifications}{
    \tightitems
        \item \textbf{SFP+ Transceivers}: Based on NIEL tests, AVAGO and CORETEK models were rejected.
        The \textbf{Direktronik 33-4248} and \textbf{FS-Europe SFP-10GSR-85} demonstrated exceptional
        resilience and were selected for final integration.
        \item \textbf{SEL Mitigation Circuit}: The supplementary over-current protection circuit was
        removed from the DB6 design. Its FDFMA3N109 MOSFETs exhibited threshold voltage shifts
        caused by the Gamma component in the reactor cavity.
        \item \textbf{DC-DC Phase Decoupling}: The LTC6902IMS QUAD oscillators were eliminated due to
        radiation-induced failures, simplifying the board while maintaining acceptable noise levels.
    \end{itemize}
    \vspace{0.12cm}
    \begin{center}
        \includegraphics[width=1.0\linewidth]{assets/fig_mosfet_threshold.png}
    \end{center}
    \vspace{-0.2cm}
    \captiontext{Threshold voltage shifts depicted in IV curves of irradiated FDFMA3N109 MOSFETs.}
}

\block{Conclusions \& Production Outlook}{
    \tightitems
        \item Radiation tests confirm that the key components of the DB6 design possess
        \textcolor{tid}{TID}, \textcolor{niel}{NIEL}, and \textcolor{see}{SEL} resistance
        well beyond the stringent HL-LHC requirements.
        \item Removing the sensitive MOSFETs and QUAD oscillators ensures the board-level design
        is highly robust.
        \item \textbf{Project Outlook}: After final design adjustments and SEU tests with operational
        firmware, the project aims to mass-produce \textbf{930 DB6 units} for the HL-LHC era,
        securing continuous and reliable data transmission for the ATLAS Tile Calorimeter.
    \end{itemize}
}

\end{columns}
\end{document}
